% Partie : counterfactuals
% Chapitre : minimal-change-semantics
% Section : contraposition

\documentclass[../../../include/open-logic-section]{subfiles}

\begin{document}

\olfileid[fr]{cnt}{min}{cpo}

\olsection{Contraposition}

Les conditionnels matériels et stricts sont équivalents à leur
contraposée. Ce n'est pas le cas des conditionnels contrefactuels.
Voici un exemple dû à Kratzer :
\begin{quote}
  Si Goethe n'était pas mort en 1832, il serait tout de même mort à ce
  jour.

  Si Goethe était encore en vie à ce jour, il serait mort en 1832.
\end{quote}
Le premier énoncé est vrai : les êtres humains ne vivent pas pendant
des siècles. Le second est manifestement faux : si Goethe était
encore en vie aujourd'hui, il ne pourrait pas être mort en 1832.

\begin{ex}\ollabel{ex:contraposition-counterex}
  La sémantique des sphères invalide la contraposition : on a $p
  \cif q \Entails/ \lnot q \cif \lnot p$. Interprétons~$p$ comme
  « Goethe n'est pas mort en 1832 » et~$q$ comme « Goethe n'est plus en
  vie aujourd'hui ». Nous pouvons représenter cette situation par un
  modèle $\mModel{M} = \tuple{W, O, V}$ où $W =
  \{w, w_1, w_2\}$, $O_w = \{\{w\}, \{w, w_1\}, \{w, w_1, w_2\}\}$,
  $V(p) = \{w_1, w_2\}$ et $V(q) = \{w, w_1\}$. % Précision facultative conservée hors du texte traduit.
%
%
   Le monde~$w$ est le monde actuel, où Goethe est mort en 1832
  et n'est plus en vie ; $w_1$ est un monde proche où Goethe est mort,
  disons, en 1833, et n'est plus en vie ; $w_2$ est un monde plus
  éloigné où Goethe est encore vivant.
\begin{figure}
\begin{center}
\begin{tikzpicture}[scale=.6,layerwidth=1.5]\tiny
  \spheresystem{3}
  \begin{scope}
    \clip \propositionplot{0}{.8}{50}{4.5} -- (4.5,-4.5) -- (-4.5,-4.5) -- (-4.5,4.5) -- (4.5,4.5);
    \spherelayer{3}
  \end{scope}
  \propositionintersect{0}{2}{110}{5.5}
  \proposition{0}{.8}{50}{4.5}
  \path (0,0) node[world] {$w$};
  \spherepos{0}{2}{node[world] {$w_1$}}
  \spherepos{-50}{3}{node[world] {$w_2$}}
  \spherepos{28}{3}{node {$q$}}
  \spherepos{42}{3}{node {$\lnot q$}}
  \spherepos{-55}{4}{node {$p$}}
  \spherepos{-67}{4}{node {$\lnot p$}}
\end{tikzpicture}
\caption{Contre-exemple à la contraposition}
\ollabel{fig:contraposition}
\end{center}
\end{figure}
Il existe une sphère admettant~$p$, $S =
\{w, w_1\}$, et $p \lif q$ est vrai dans tous ses mondes ; donc
$\mSat{M}{p \cif q}[w]$. En revanche, la sphère admettant~$\lnot q$,
$\{w, w_1, w_2\}$, contient le monde~$w_2$, où
$q$ est faux et~$p$ est vrai ; donc
$\mSat/{M}{\lnot q \lif \lnot p}[w_2]$. Le conditionnel obtenu par
contraposition est ainsi faux au monde central.
\end{ex}

\end{document}
